\chapter{Computational Number Representations}
\label{chap:numbers}

For our computational geometry algorithms to be correct, we need to use exact
number representations --- roundoff error introduced by floating-point
approximations can easily lead to incorrect outcomes~\cite{kettner2008}.
For example, if the intersection of two polygons is a single point or line, small rounding
errors could easily make the intersection look empty.

Mathlib provides an exact real number representation that encodes a number as
an infinite sequence of rational approximations. Unfortunately, this
representation is not suitable for computation because some operations are not
computable. For example, determining whether two reals are equal requires
iterating over the entire infinite sequences that represent them.

The rational numbers ($ℚ$) are a reasonable choice of number type: fractions are easy
to compute with and don't introduce any rounding error. Some applications
require irrational numbers though. For example, rotating a unit line segment by
$45°$ produces a point with coordinates $(\sqrt{2}/2, \sqrt{2}/2)$, whose
coordinates are irrational.

If we only cared about $45°$ rotations, we could store numbers of the
form $a_1 + a_2\sqrt{2}$ where $a_1$ and $a_2$ are rational numbers. It
turns out that adding, multiplying, and dividing numbers of this form produces
numbers of this form, and it is not too hard to compare them for equality and
inequalities; these operations are sufficient to perform many geometric
operations (including polygon union and intersection).

Of course if you give a mouse a $45°$ rotation, they will want a
rotation of $30°$. We can support both $45°$ and $30°$
rotations with numbers of the form
$a_1 + a_2\sqrt{2} + a_3\sqrt{3} + a_6\sqrt{6}$; these numbers also support the
important operations for many geometric algorithms.

This suggests that the choice of number representation should be application
specific, and the geometric algorithms should be compatible with different
number representations. This package constructs a few concrete number
representations that applications can use, and proves that they are correct.

In this context, correctness means that the numbers support the necessary
operations, and that those operations behave as you would expect. More
concretely, we show that the number representations form ``linearly ordered
fields,'' which are types that support addition, multiplication, subtraction,
division, and computable comparisons (equality and inequality), all of which
must satisfy certain axioms. The geometric constructions in the remainder of
the development are designed to work with any linearly ordered field.

The first implementation is \code{AdjoinSqrt A n} (abbreviated $A[\sqrt{n}]$),
which represents numbers of the form $a_1 + a_n\sqrt{n}$. Instead of requiring
$a_1$ and $a_n$ to be rationals, we allow them to be taken from any fixed
linearly ordered field $A$.  This allows the construction to be repeated. For
example, we can factor numbers of the form $a_1 + a_2\sqrt{2} + a_3\sqrt{3} +
a_6\sqrt{6}$ as $(a_1 + a_2\sqrt{2}) + (a_3 + a_6\sqrt{2})\sqrt{3}$, and then
represent them using $ℚ[\sqrt{2}][\sqrt{3}]$.

The second implementation is intended to improve performance by using
fast approximations wherever possible, and only falling back on the
exact computations when the approximations are not precise enough. We define
\code{Filtered $A$} which stores both an unevaluated exact expression for a
value and a fast approximation of the value; it keeps the approximations up to
date and only does the slower computation when the approximations aren't
sufficient to determine the result of a comparison.

In both implementations, we prove that if the underlying type $A$ is a linearly
ordered field, then the constructed type \code{AdjoinSqrt $A$ $n$} or
\code{Filtered $A$} is also a linearly ordered field.
Section~\ref{sec:numberType} contains the little details necessary to prove that
\code{Filtered $ℚ[\sqrt{2}][\sqrt{3}]$} is a linearly ordered field.

The definition of a linearly ordered field can be broken down into parts, and
these parts are important stepping stones on the path to proving that a number
representation is a linearly ordered field. In particular, we will mention the
following intermediate interfaces:

\begin{enumerate}
\item A (commutative) ring must support addition, subtraction, and
multiplication, but not division. For example, the integers ($ℤ$) are a ring
but not a field. Rings are useful in this context because one can generate a
field from a ring by forming fractions; an alternative approach to defining
$ℚ[\sqrt{2}][\sqrt{3}]$ would be to take the field of fractions of
$ℤ[\sqrt{2}][\sqrt{3}]$.

\item A field is a ring that also supports division.

\item A linear order is a type that supports comparisons ($≤$, $<$, etc).

\item A linearly ordered ring is a ring with a linear order that is compatible
with the ring operations; a linearly ordered field is a field with a linear
order that is compatible with the field operations.
\end{enumerate}

Note that we are borrowing all of these definitions from Mathlib, so we are
leaving the exact definitions a bit vague here.  These definitions can also be
broken down further (e.g. Monoid, Group, Partial Order, etc) and we do prove
related intermediate results where possible, but these are probably less useful
for this library so we don't describe them here.

\section{Adjoining a square root}

%% Definition

\begin{definition}[Adjoining a square root]
  \label{def:adjoin}
  \lean{AdjoinSqrt}
  \leanok
  An element of $A[\sqrt{n}]$ is stored as a pair $(a_1, a_n)$, representing the value
  $a_1 + a_n\sqrt{n}$.
\end{definition}

We begin by checking that $A[\sqrt{n}]$ is a ring. We will then proceed to check
that it is a field, and that it has a natural ordering on it ($<$ and $≤$).

\begin{definition}[{Ring operations on $A[\sqrt{n}]$}]
  \label{def:adjoinRingOps}
  \uses{def:adjoin}
  \lean{AdjoinSqrt.instRingOps}
  \leanok
  We define the basic ring operations of addition, negation, and multiplication
  in the usual way.
  For example, $(a_1 + a_n\sqrt{n})(b_1 + b_n\sqrt{n}) = (a_1b_1 + na_nb_n) + (a_1b_n + a_nb_1)\sqrt{n}$.

  Note that this definition only makes sense if $A$ supports these operations.
\end{definition}

%% A[√n] is a ring

\begin{proposition}
  \label{thm:adjoinRing}
  \lean{AdjoinSqrt.instCommRing}
  \leanok
  \uses{def:adjoinRingOps}
  If $A$ satisfies the ring properties, then so does $A[\sqrt{n}]$.
\end{proposition}

\begin{proof}
  \leanok
  Follows immediately from the definitions.
\end{proof}

%% A[√n] has √n

Of course, the whole point of $A[\sqrt{n}]$ is that it contains a square root of $n$:
\begin{theorem}
  \label{thm:adjoinHasSqrtN}
  \lean{AdjoinSqrt.root_mul_self}
  \uses{def:adjoin, def:adjoinOfA, thm:adjoinRing}
  \leanok
  If $A$ is a commutative ring, then in $A[\sqrt{n}]$, $\sqrt{n} \cdot \sqrt{n} = n$.
\end{theorem}

\begin{proof}
  \leanok
  Follows immediately from the definitions.
\end{proof}

%% A[√n] is a field

In order to check that $A[\sqrt{n}]$ is a field, we need to define division. In
fact, we only need to define $x^{-1}$ for $x ∈ A[\sqrt{n}]$, which we do as
follows:

\begin{definition}[Inverse]
  \label{def:adjoinInv}
  \lean{AdjoinSqrt.instInv}
  \leanok
  \uses{def:adjoinRingOps, def:norm}
  We define $(a_1 + a_n\sqrt{n})^{-1}$ as follows:
  \begin{eqnarray*}
    \frac{1}{a_1 + a_n\sqrt{n}}
      & = & \frac{1}{a_1 + a_n\sqrt{n}} \cdot \frac{a_1 - a_n\sqrt{n}}{a_1 - a_n\sqrt{n}}
        =   \frac{a_1 - a_n\sqrt{n}}{a_1^2 - na_n^2} \\ \\
      & = & a_1(a_1^2 - na_n^2)^{-1} - a_n(a_1^2 - na_n^2)^{-1}\sqrt{n} \\
  \end{eqnarray*}
  As long as $A$ has inverses, this is an element of $A[\sqrt{n}]$.
\end{definition}

\begin{lemma}
  \label{thm:adjoinField}
  \lean{AdjoinSqrt.instField}
  \leanok
  \uses{def:adjoin, thm:adjoinRing, def:adjoinInv, def:norm}
  If $A$ is a field and $n$ is not a square in $A$, then $A[\sqrt{n}]$ is a
  field.
\end{lemma}

\begin{proof}
  \leanok
  This basically boils down to the computation given in the definition above.

  However, we need the extra condition that $n$ is not a square because
  otherwise $a_1^2 - na_n^2$ might be 0 even if $a$ isn't zero, in which case
  $a$ wouldn't have an inverse.
\end{proof}

%% Norm

Two expressions in the definition above are quite useful in the proofs because they allow relating
values in $A[\sqrt{n}]$ with values in $A$:

\begin{definition}[Conjugate]
  \label{def:conj}
  \lean{AdjoinSqrt.conj}
  \uses{def:adjoinRingOps}
  \leanok
  We refer to $a_1 - a_n\sqrt{n}$ as the \emph{conjugate} of $a$ (written
  $\conj{a}$).
\end{definition}
\begin{definition}[Norm]
  \label{def:norm}
  \lean{AdjoinSqrt.norm}
  \uses{def:conj}
  \leanok
  We refer to $a_1^2 - na_n^2$ as the \emph{norm} of $a$ (written $\norm{a}$).
  Note that it is an element of $A$, not just $A[\sqrt{n}]$. With these
  definitions, we have $\norm{a} = a \cdot \conj{a}$ and $a^{-1} =
  \conj{a}\norm{a}^{-1}$.
\end{definition}

%% A[√n] is ordered ring

Now we can proceed to check the ordering properties on $A[\sqrt{n}]$. Instead
of defining comparisons like $<$ and $≤$ directly, it is easier to start by
defining a comparison with 0. Then we can define $a \leq b$ if and only if $0
\leq b - a$. The set of values greater than or equal to 0 is referred to as a
\emph{positive cone}.

As with the other operations, we will use an ordering on $A$ to define the
positive cone on $A[\sqrt{n}]$. The key idea is that we compare the squares
to see if $|a_1|$ or $|a_n\sqrt{n}|$ is larger, and then take the sign of
whichever one dominates.
\begin{definition}[{Positive cone on $A[\sqrt{n}]$}]
  \label{def:posCone}
  \lean{AdjoinSqrt.IsNonneg}
  \leanok
  \uses{def:adjoin}
  We say $0 \leq a_1 + a_n\sqrt{n}$ if either $a_1^2 \geq na_n^2$ and $a_1 \geq
  0$ (i.e. if $a_1$ dominates and is nonnegative) or if $a_1^2 \leq na_n^2$ and $a_n \geq 0$
  (i.e. if $a_n\sqrt{n}$ dominates and is nonnegative).

  Note that we can rewrite $a_1^2 \geq na_n^2$ as $\norm{a} \geq 0$.
\end{definition}

In order to use the positive cone $0 \leq \cdot$ to define a linearly ordered
ring, we need to check that the positive cone interacts properly with the ring
operations. The full set of axioms and the proof that a
well-behaved positive cone can be turned into a linearly ordered ring are
defined by Mathlib:

\begin{definition}[Well-behaved positive cone]
  \label{def:ringCone}
  \lean{RingCone, IsOrderedRing.mkOfCone}
  \mathlibok
  A well-behaved positive cone on a ring $A$ is a relation $0 \leq \cdot$ that satisfies:
  \begin{itemize}
    \item $0 \leq 0$
    \item $0 \leq 1$
    \item either $0 \leq a$ or $0 \leq -a$
    \item (antisymmetry) if $0 \leq a$ and $0 \leq -a$ then $a = 0$
    \item (multiplicativity) if $0 \leq a$ and $0 \leq b$ then $0 \leq ab$
    \item (additivity) if $0 \leq a$ and $0 \leq b$ then $0 \leq a + b$
  \end{itemize}
\end{definition}


What's left for us is to prove that our positive cone definition for
$A[\sqrt{n}]$ is well-behaved. The first three properties follow immediately
from the definitions and the fact that $A$'s order is well-behaved; the latter
three are more tricky and we state them as lemmas:

\begin{lemma}[Antisymmetry]
  \label{thm:antisymm}
  \lean{AdjoinSqrt.nonneg_antisymm}
  \leanok
  \uses{def:posCone}
  If $A$ is a linearly ordered \emph{field} and $n$ is a positive non-square in
  $A$ then $A[\sqrt{n}]$ satisfies the antisymmetry property: if $0 \leq a$ and
  $0 \leq -a$ then $a = 0$.
\end{lemma}
\begin{proof}
  \leanok
  If $a$ is $a_1$-dominated then so is $-a$, and we can conclude that $a_1 =
  0$, so $a_n$ must also be 0, so $a = 0$. Similarly, if $a$ is $a_n$-dominated
  then $a$ must be 0.

  The only tricky corner case is if $\norm{a} = 0$. This case requires that $A$
  is a field and that $n$ is non-square.  But we have $a_1^2 = na_n^2$ so that
  $n = (a_1/a_n)^2$, a contradiction unless $a_1 = a_n = 0$.
\end{proof}

\begin{lemma}[Multiplicativity]
  \label{thm:multiplicativity}
  \lean{AdjoinSqrt.IsNonneg.mul}
  \leanok
  \uses{def:posCone}
  If $A$ is a linearly ordered ring and $n$ is positive in $A$, then
  $A[\sqrt{n}]$ satisfies the multiplicative property: for all $a, b ∈
  A[\sqrt{n}]$, if $0 \leq a$ and $0 \leq b$ then $0 \leq ab$.
\end{lemma}
\begin{proof}
  \leanok
  You can split into four cases based on whether $a$ is $a_1$-dominated and
  whether $b$ is $b_1$-dominated; since $\norm{ab} = \norm{a}\norm{b}$ you can
  determine the sign of $\norm{ab}$ in each of the cases which tells you which
  coefficient of $ab$ you need to prove is nonnegative.

  In each of the cases, the coefficient in question is a sum of two terms, one
  of which is obviously nonnegative from the assumptions. The goal is then to
  prove that the other term is smaller in absolute value (so that the sum is
  also nonnegative), which you can do by writing out the difference of the squares.

  For example, let's take the case where $a$ is $a_1$-dominated and $b$ is
  $b_n$-dominated. Then we know $\norm{ab} = \norm{a}\norm{b} \leq 0$ (so that
  $ab$ is $(ab)_n$-dominated), $a_1 \geq 0$, and $b_n \geq 0$. To show $ab \geq
  0$ we must show that $(ab)_n \geq 0$. We have $(ab)_n = a_1b_n + b_1a_n$; the
  first term is nonnegative so we just need to show that $|b_1a_n| \leq |a_1b_n|$.

  Comparing the squares, we have
  \begin{displaymath}
    (a_1b_n)^2 - (b_1a_n)^2 = b_n^2\norm{a} - a_n^2\norm{b}
  \end{displaymath}
  the two terms on the right are both nonnegative since $\norm{a} \geq 0$ and
  $\norm{b} \leq 0$. Therefore $|b_1a_n| \leq |a_1b_n|$, so we have $(ab)_n
  \geq 0$ as required.

  The other three cases are almost identical, although the cases where both
  norms have the same sign also require $n > 0$ for the inequalities to work out.
\end{proof}

\begin{lemma}[Additivity]
  \label{thm:additivity}
  \lean{AdjoinSqrt.IsNonneg.add}
  \leanok
  \uses{def:posCone}
  If $A$ is a linearly ordered \emph{field} and $n$ is a positive non-square in $A$, then
  $A[\sqrt{n}]$ satisfies the additive property: for all $a, b ∈
  A[\sqrt{n}]$, if $0 \leq a$ and $0 \leq b$ then $0 \leq a + b$.
\end{lemma}
\begin{proof}
  \leanok
  \uses{thm:multiplicativity, thm:antisymm}
  This is the trickiest proof of this section. To keep track of the cases, it
  is helpful to organize them into the following diagram:

  \begin{center}
    \begin{tikzpicture}[scale=1.2]
      % Axes: a₁ horizontal, aₙ√n vertical
      \draw[<->] (-3,0) -- (3,0) node[right] {$a_1$};
      \draw[<->] (0,-3) -- (0,3) node[above] {$a_n\sqrt{n}$};

      % |a₁| = |aₙ√n|: separates a₁-dominated from aₙ-dominated
      \draw[dashed] (-2.5,-2.5) -- (2.5,2.5)
        node[pos=0.9, sloped, above, align=center, font=\small]
          {$\norm{a} \leq 0$ \\ $a_n$-dominated}
        node[pos=0.9, sloped, below, align=center, font=\small]
          {$a_1$-dominated \\ $\norm{a} \geq 0$};

      % a₁ + aₙ√n = 0: separates a ≥ 0 from a ≤ 0
      \draw[dashed] (-2.5,2.5) -- (2.5,-2.5)
        node[pos=0.85, sloped, above, font=\small] {$a \geq 0$}
        node[pos=0.85, sloped, below, font=\small] {$a \leq 0$};

      % Quadrant names
      \node at (0,3.8) {N};
      \node at (3.8,0) {E};
      \node at (0,-3.8) {S};
      \node at (-3.8,0) {W};
    \end{tikzpicture}
  \end{center}

  I'll refer to quadrants of the diagram by their cardinal directions; for
  example $a$ is in the north quadrant (N) when $\norm{a} \leq 0$ and $a \geq 0$,
  while $a$ is in the NE quadrant if $a_1 \geq 0$ and $a_n \geq 0$. Our
  assumption that $a$ and $b$ are nonnegative says that $a$ and $b$ live in the
  NE half-plane. Anything in the NE quadrant is obviously nonnegative because
  regardless of which coefficient dominates, it is nonnegative.

  As with multiplicativity, you can divide into cases depending on whether $a$
  and $b$ are first- or second-coefficient dominated. The difficulty arises
  because the goal $0 \leq a + b$ is a statement about the coefficients $a_i$
  and $b_i$, but the assumptions about the norm talk about $a_i^2$ and $b_i^2$.

  The trick is to add a multiplier so that the goal becomes a statement about
  the squares of $a_i$ and $b_i$, which makes the assumptions about the norms
  more useful. Instead of proving $0 \leq (a + b)$ directly, we instead prove
  that $0 \leq (a + b)m$ for a well-chosen and obviously-nonnegative $m$, and then
  use the multiplicative property to conclude that $0 \leq a + b$.

  To choose $m$, we need something with terms that come from $a$ and $b$ or
  their conjugates, and we need something obviously nonnegative. It turns out that
  summing the absolute values of the coefficients works out: it is obviously
  nonnegative because both $|a_1|$ and $|a_n|$ are nonnegative so either branch of
  the definition of nonnegativity gives a nonnegative value. It is useful because
  $|a_1| + |a_n|\sqrt{n}$ is always either $a$, $\conj{a}$, $-a$, or
  $-\conj{a}$, so many of the terms of $(a + b)m$ are either $\pm a^2$ or $\pm
  a\conj{a} = \pm \norm{a}$.

  Specifically, we choose $m = (|a_1| + |a_n|\sqrt{n}) + (|b_1| + |b_n|\sqrt{n})$.

  In the diagram, everything in the NE quadrant is obviously nonnegative, and
  we choose $m$ by reflecting $a$ and $b$ into the NE quadrant.

  Even after this transformation there are still a bunch of cases that require
  some analysis. In each case, we expand $(a+b)m$ and show that it is obviously
  nonnegative. We'll prove two of them here; the remainder are similar.

  First, consider the case where $a$ is in the NNW wedge (i.e. $a_1 \leq 0$ and
  $a \geq 0$), and $b$ is in the ESE wedge. In this case, we have $m = \conj{b}
  - \conj{a}$, so
  \begin{eqnarray*}
    (a + b)m
      & = & b\conj{b} - a\conj{a} + a\conj{b} - b\conj{a} \\
      & = & \norm{b} - \norm{a} + 2(a_nb_1 - a_1b_n)\sqrt{n} \\
  \end{eqnarray*}

  Now, because of the wedges that $a$ and $b$ are in, we have $\norm{a} \leq 0
  \leq \norm{b}$ so the $A$ coefficient $\norm{b} - \norm{a} \geq 0$ (we can
  use additivity here because the norms are in $A$, not $A[\sqrt{n}]$).
  We also have $0 \leq a_1^2 \leq na_n^2$ and $0 \leq nb_n^2 \leq b_1^2$;
  multiplying these we get $na_1^2b_n^2 \leq na_n^2b_1^2$. Since $n$ is
  positive, we can conclude $a_1b_n \leq |a_1b_n| \leq |b_1a_n| = a_nb_1$ so the
  $\sqrt{n}$ coefficient $2(a_nb_1 - a_1b_n)$ is also nonnegative. That means
  $(a+b)m$ is in the NE quadrant, and is therefore obviously nonnegative so
  we're done.

  Now, consider the case where both $a$ and $b$ are in the ESE wedge. In this
  case, $m = \conj{a} + \conj{b}$, so
  \begin{eqnarray*}
    (a + b)m
      & = & b\conj{b} + a\conj{a} + a\conj{b} + b\conj{a} \\
      & = & \norm{b} + \norm{a} + 2(a_1b_1 - na_nb_n) \\
  \end{eqnarray*}

  Since $a$ and $b$ are both in the E quadrant, the norm terms are nonnegative.
  We also have $0 \leq na_n^2 \leq a_1^2$ and $0 \leq nb_n^2 \leq b_1^2$;
  multiplying these gives $n^2a_n^2b_n^2 \leq a_1^2b_1^2$ (and we know $a_1b_1
  \geq 0$), so the sum of the other two terms is nonnegative. As in the
  previous case, $(a + b)m$ is obviously nonnegative.

  The remaining cases are similar. In all cases, we conclude that $m$ and $(a +
  b)m$ are obviously nonnegative. If $a + b$ were negative, then $m$ would have
  to be 0, which means all the coefficients of $a$ and $b$ are 0; this would be
  a contradiction, so we conclude $a + b$ must be nonnegative, as required.
\end{proof}

\begin{corollary}
  \label{thm:adjoinOrderedRing}
  \lean{AdjoinSqrt.nonnegCone, AdjoinSqrt.instLinearOrder, AdjoinSqrt.instIsStrictOrderedRing}
  \leanok
  If $A$ is a linearly ordered field and $n$ is a positive non-square in $A$, then
  $A[\sqrt{n}]$ is also a linearly ordered ring.
\end{corollary}
\begin{proof}
  \leanok
  \uses{thm:multiplicativity, thm:additivity, def:ringCone}
  Apply Mathlib's construction (Definition~\ref{def:ringCone}) to
  Lemmas~\ref{thm:multiplicativity} and~\ref{thm:additivity}.
\end{proof}

%% A[√n] is an ordered field

Putting this all together gives the main result of this section:

\begin{theorem}
  \label{thm:adjoinOrderedField}
  \lean{AdjoinSqrt.instField, AdjoinSqrt.instLinearOrder, AdjoinSqrt.instIsStrictOrderedRing}
  \uses{def:adjoin, thm:adjoinField, thm:adjoinOrderedRing}
  \leanok

  If $A$ is a linearly ordered field, and if $n$ is a positive non-square in $A$
  then $A[\sqrt{n}]$ is also a linearly ordered field.
\end{theorem}

\begin{proof}
  \leanok
  The field structure is Lemma~\ref{thm:adjoinField} and the order is
  Corollary~\ref{thm:adjoinOrderedRing}; since both are instances on the same
  type, nothing further is needed to combine them.
\end{proof}

%%%%%%%%%%%%%%%%%%%%%%%%%%%%%%%%%%%%%%%%%%%%%%%%%%%%%%%%%%%%%%%%%%%%%%%%%%%%%%%%
\subsection{Converting from $A$ to $A[\sqrt{n}]$}

These are some results that aren't directly needed to use these number representations, but are
interesting nonetheless.

%% A → A[√n]

\begin{definition}[Embedding of $A$]
  \label{def:adjoinOfA}
  \lean{AdjoinSqrt.instCoe}
  \leanok
  \uses{def:adjoin}
  We write $A \hookrightarrow A[\sqrt{n}]$ for the natural embedding
  $a \mapsto a + 0\sqrt{n}$ (assuming $0 ∈ A$).
\end{definition}

%% A[√n] is an A-algebra

\begin{lemma}
  \label{thm:adjoinAlgebra}
  \lean{AdjoinSqrt.algebra}
  \leanok
  \uses{def:adjoin, def:adjoinOfA, thm:adjoinRing}
  If $A$ is a ring, then $A[\sqrt{n}]$ is an $A$-algebra.
\end{lemma}

\begin{proof}
  \leanok
  Trivial
\end{proof}

\begin{definition}[Quadratic algebra]
  \label{def:quadraticAlgebra}
  \lean{QuadraticAlgebra}
  \mathlibok
  For a commutative semiring $A$ and $a, b ∈ A$, the quadratic algebra over $A$
  with coefficients $a$, $b$ has elements $x + yω$ for $x, y ∈ A$, with
  $ω^2 = a + bω$.
\end{definition}

\begin{lemma}
  \label{thm:mathlibAdjoin}
  \lean{AdjoinSqrt.quadraticAlgebraEquiv}
  \leanok
  \uses{def:adjoin, thm:adjoinAlgebra, def:quadraticAlgebra}
  If $A$ is a commutative semiring, then $A[\sqrt{n}]$ is isomorphic, as an
  $A$-algebra, to the quadratic algebra over $A$ with coefficients $n$ and $0$,
  by $a_1 + a_n\sqrt{n} \mapsto a_1 + a_nω$.
\end{lemma}

\begin{proof}
  \leanok
  Trivial.
\end{proof}

\TODO{Talk about sign-based formulation and how it avoids the addition and
multiplication in the cone computation. Link it to the implementation, cite
lemma that two definitions are equivalent.}

%%%%%%%%%%%%%%%%%%%%%%%%%%%%%%%%%%%%%%%%%%%%%%%%%%%%%%%%%%%%%%%%%%%%%%%%%%%%%%%%
\section{Approximate arithmetic}
Rational numbers are exact, but they can be quite expensive to compute with
because the size of the numerator and denominator can grow unboundedly. You can
reduce the storage needed by always reducing common factors between the
numerator and denominator, but then you need to perform a gcd operation on every
arithmetic operation.

There are fast approximate ways to represent numbers, such as floating point
numbers or dyadic rationals (numbers of the form $m/2^k$ with integers $m$ and
$k$). By keeping the size of the number representations bounded, we store and
compute with these numbers efficiently, but we give up precision.

The \code{Filtered} type defined in this section takes the best of both
worlds, much as CGAL's filtered kernels do~\cite{cgal}. We store a computation
that evaluates to the exact value, but we keep it unevaluated. We also store a fast approximation of that value that we update on
every operation. We only fall back to the exact computation if the approximation
isn't sufficient to compare two values.

To define our \code{Filtered} type, we first need to formalize what we need
from an approximation strategy. First, we need an approximation type that
supports fast versions of the operations we care about. The intuition behind the
definition is that the approximations are fuzzy, so that an approximate value
$X$ could potentially represent many real numbers.
\begin{definition}[Approximation type]
  \label{def:realApprox}
  \lean{CGLean.RealApprox}
  \leanok
  An approximation type $α$ is a type whose values represent subsets of the real
  numbers, and that supports the ring operations. The ring operations must be
  compatible with the operations on the approximated real numbers (e.g. if $x ∈
  X$ and $y ∈ Y$ then $x + y ∈ X + Y$, and similarly for
  multiplication and subtraction).

  An inverse operation should also be defined, but it may not be possible to
  represent the inverse of an approximate value, especially if that value could
  represent 0. For example, there is no interval that contains the inverses of
  all possible numbers between $-1$ and $1$. We do require that if the inverse
  exists, then it is compatible with the inverse on the real numbers in the same
  way that we require for the ring operations.

  Similarly, the approximation should have a partial comparison operation, but
  we may not be able to determine the order between two approximations if they
  overlap. However, we do require that if $X < Y$ then for all $x ∈
  X$ and $y ∈ Y$, that $x < y$, and similarly for $=$ and $>$.
\end{definition}
Notably, we \emph{don't} require the ring properties for approximation types; in
fact we don't really need a notion of equality.

This definition doesn't require that the partial inverse and comparison
operations \emph{only} fail when they must (i.e. when trying to invert an
approximation that could be 0, or trying to compare approximations that
overlap). However, those can be nice properties to check for an implementation,
so we also have a definition that includes those requirements:

\begin{definition}[Complete approximation type]
  \label{def:completeApprox}
  \lean{CGLean.RealApprox.IsComplete}
  \uses{def:realApprox}
  \leanok
  An approximation type is \emph{complete} if \begin{itemize}
    \item For all $X ∈ α$, if $0 \notin X$ then $X^{-1}$ is defined
    \item If the comparison between $X$ and $Y$ is undefined, then there is
      some $x_1 ∈ X, y_1 ∈ Y$ and some $x_2 ∈ X, y_2 ∈ Y$ such that
      comparing $x_1$ to $y_1$ gives a different answer than comparing $x_2$ to
      $y_2$.
  \end{itemize}
\end{definition}

We have implemented the approximation type using intervals whose endpoints are
bounded-precision dyadic rationals; see Subsection~\ref{sec:dyadicInterval}. The
\code{Filtered} type is generic though, so you can also implement an
approximation type using floating-point arithmetic, fixed-point arithmetic
stored in machine words, or a variety of other formats.

Once you have a fast approximation type, you also need a strategy $f : A → α$
for approximating exact values in $A$ by approximate values in $α$. In order to
be useful, $f$ must be compatible with the embeddings of $A$ and $α$ into $ℝ$:

\begin{definition}[Valid approximation strategy]
  \label{def:approxStrategy}
  \uses{def:realApprox}
  \lean{CGLean.ApproxStrategy}
  \leanok
  An approximation strategy for a linearly ordered ring $A$ is an embedding $i :
  A → ℝ$ (i.e. a function $i$ that preserves the linearly ordered ring
  operations), an approximation type $α$, and a function $f : A → α$ such that
  for all $a ∈ A$, $i(a) ∈ f(a)$. We typically just refer to $f$, since $i$ is
  unique and usually clear from context.
\end{definition}

Given a valid approximation strategy, we can now define the \code{Filtered}
type:

\begin{definition}[Filtered number representation]
  \label{def:filtered}
  \uses{def:approxStrategy}
  \lean{CGLean.Filtered, CGLean.Filtered.Pre, CGLean.Filtered.get}
  \leanok
  Given an approximation strategy $f : A → α$, a \code{Filtered $A$ $f$}
  value consists of an unevaluated value $a$ of type $A$ and an approximation $X
  ∈ α$ with $i(a) ∈ X$. We consider two values in \code{Filtered $A$ $f$} to be
  equal if the wrapped $A$ values are equal.
\end{definition}

\begin{definition}[Ring operations]
  \label{def:filteredRingOps}
  \uses{def:filtered}
  \lean{CGLean.Filtered.instRingOps}
  \leanok
  We define the usual ring operations on \code{Filtered} values by storing
  thunks that perform the corresponding operation on the underlying $A$ value.
  We also immediately compute a new approximation using the corresponding operations
  on $α$.

  For this definition to make sense, we need to ensure that the uncomputed
  exact values and the immediately computed approximate values match up (e.g.
  that $i(a + b) ∈ X + Y$). This follows because $i$ is an
  embedding, so that $i(a + b) = i(a) + i(b)$, and because $α$'s
  operations are compatible with the real number operations, so that $i(a) +
  i(b) ∈ X + Y$. The same reasoning applies for the other operations.
\end{definition}

\begin{lemma}
  \label{thm:filteredRing}
  \lean{CGLean.Filtered.instCommRing}
  \uses{def:filteredRingOps}
  \leanok
  If $A$ is a commutative ring, then \code{Filtered $A$ $f$} is a commutative
  ring.
\end{lemma}

\begin{proof}
  \leanok
  Since there's a 1-1 correspondence between elements of $A$ and elements of
  \code{Filtered $A$ $f$} that preserves all of the operations, this follows
  immediately from the corresponding property for $A$.
\end{proof}

\begin{definition}[Filtered inverse]
  \label{def:filteredInv}
  \uses{def:filtered}
  \lean{CGLean.Filtered.Pre.inv, CGLean.Filtered.instInv}
  \leanok
  If $A$ has an inverse operation, we define an inverse on \code{Filtered $A$
  $f$} by taking an inverse of the approximate value if it has an inverse, and
  otherwise computing the exact inverse in $A$ and reapproximating it.
\end{definition}

\begin{lemma}
  \label{thm:filteredField}
  \lean{CGLean.Filtered.instField}
  \uses{def:filteredInv, thm:filteredRing}
  \leanok
  If $A$ is a field, then \code{Filtered $A$ $f$} is a field.
\end{lemma}

\begin{proof}
  \uses{thm:filteredRing}
  \leanok
  Follows immediately.
\end{proof}

\begin{definition}[Filtered comparison]
  \label{def:filteredCompare}
  \uses{def:filtered}
  \lean{CGLean.Filtered.Pre.cmp, CGLean.Filtered.Pre.cmp_eq, CGLean.Filtered.compare_get}
  \leanok
  We define the ordering on \code{Filtered $A$ $f$} using the approximate
  comparison and only fall back to the ordering on $A$ if the approximation is
  inconclusive.

  The requirement that the approximation strategy is valid and that the
  approximation type's comparison function is compatible with the comparison on
  the real numbers means that regardless of whether we take the fast or the slow
  path, we still get the same answer as we would have in $A$.
\end{definition}

\begin{lemma}
  \label{thm:filteredOrdered}
  \lean{CGLean.Filtered.instLinearOrder, CGLean.Filtered.instIsStrictOrderedRing}
  \uses{thm:filteredRing}
  \leanok
  If $A$ is a linearly ordered commutative ring, then so is \code{Filtered $A$
  $f$}, and its comparison agrees with that of $A$.
\end{lemma}

\begin{proof}
  \uses{thm:filteredRing, def:filteredCompare}
  \leanok
  Follows immediately from the fact the comparison matches that on $A$.
\end{proof}

\subsection{Approximating with dyadic intervals}
\label{sec:dyadicInterval}
While the \code{Filtered} type works with many approximation strategies, we have
built a simple approximation type based on dyadic rationals as a proof of
concept. The definitions of dyadic intervals and exact operations on them are
taken from the LeanCert library:

\begin{definition}[Dyadic interval]
  \label{def:dyadicInterval}
  \lean{LeanCert.Core.IntervalDyadic}
  \mathlibok % It's not in mathlib, but it isn't here either
  A \emph{dyadic interval} is a closed interval whose endpoints are dyadic
  rationals: numbers of the form $m\cdot2^k$ for integers $m$ and $k$.
\end{definition}

For this development, we define all of the operations on dyadic intervals so
that they round the intervals out to a fixed precision:

\begin{definition}[Operations on dyadic intervals]
  \label{def:dyadicOps}
  \lean{CGLean.RoundedInterval.trim, CGLean.RoundedInterval.compare?, CGLean.RoundedInterval.inv?}
  \uses{def:dyadicInterval}
  \leanok
  We define the ring operations on our dyadic intervals by first performing the
  exact ring operations on the endpoints, and then rounding them in the correct
  direction to dyadic rationals with a fixed precision (i.e. numbers of the form
  $m \cdot 2^k$ with $m$ having at most 64 bits).

  We define inversion the same way, but only if the interval doesn't contain 0.

  We define comparison by comparing the upper end of one interval with the
  other: when comparing the intervals $X = [x_{lo}, x_{hi}]$ and $Y = [y_{lo},
  y_{hi}]$:\begin{itemize}
    \item If $x_{hi} < y_{lo}$ then we say $X < Y$
    \item If $x_{lo} > y_{hi}$ then we say $X > Y$
    \item If $x_{lo} = x_{hi} = y_{lo} = y_{hi}$ then we say $X = Y$
    \item Otherwise the comparison is inconclusive
  \end{itemize}
\end{definition}

\begin{lemma}[Dyadic intervals are an approximation type]
  \label{thm:dyadicApprox}
  \lean{CGLean.RoundedInterval.instDyadicApprox}
  \uses{def:dyadicOps, def:realApprox}
  \leanok
  Dyadic intervals are an approximation type, i.e. the operations on the
  approximations produce approximations of the corresponding operations on
  real numbers.
\end{lemma}

\begin{proof}
  \leanok
  This boils down to checking that we're rounding in the right direction and the
  basic properties about dyadic arithmetic that LeanCert provides.
\end{proof}

We can also check that we only fail to invert or compare if we need to:
\begin{lemma}
  \label{thm:dyadicComplete}
  \lean{CGLean.RoundedInterval.instDyadicComplete}
  \uses{def:completeApprox, thm:dyadicApprox}
  \leanok
  Dyadic intervals are a complete approximation type.
\end{lemma}

\begin{proof}
  \leanok
  We need to show that comparison is only undefined if the intervals overlap and
  that inversion is undefined only if the interval contains zero. The second
  property is obvious from the definition. To check the comparison property, if
  none of the comparison definitions apply then the intervals overlap; you can
  use the endpoints of the intervals themselves as the examples of numbers that
  don't compare consistently.
\end{proof}

The dyadic interval approximation is probably not the most efficient choice for
our fast approximation. Even though we round to 64 bits, Lean's fast small
integer implementation only applies at 31 bits, and all of the dyadic operations
blow up the precision before rounding it back down.

Nevertheless, a simple benchmark shows that even the unoptimized dyadic filter
provides an improvement (see \code{bench/Main.lean}):

\begin{center}
\begin{tabular}{lrrrr}
  \hline
  orientation tests in $ℚ[\sqrt{2}][\sqrt{3}]$ & tests & exact & filtered & fallbacks \\
  \hline
  random input points & 2000 & 88 ms & 20 ms & 21 \\
  intersections of input lines & 2000 & 651 ms & 39 ms & 20 \\
  intersections of lines through those & 2000 & 2815 ms & 134 ms & 55 \\
  all collinear & 100 & 38 ms & 45 ms & 100 \\
  \hline
\end{tabular}
\end{center}
The intersection operation produces coordinates by algebraically operating on
the input points, which has the effect of increasing the size of the exact
representations. That's why the runtime increases in each of the first three
rows. You can see that the filtering reduces a lot of computational overhead in
these cases.

The ``fallbacks'' column gives the number of times that the filter was
insufficient and the exact arithmetic was needed. In the ``all collinear'' case,
all of the tested computations produce equality, which means the filter doesn't
help at all. This row shows that the overhead of doing the filtering is small
($\sim$18\%).

It's also worth noting that for this test, the fallbacks were only taken when
the numbers being compared were actually equal, which means that the filter was
always useful on random data.

A floating-point implementation would almost certainly be an even bigger
improvement, but would require axiomatizing the IEEE 754 floating point
specification (similar to how Flocq~\cite{flocq} does for Coq). We leave that
formalization to future work.

Degenerate cases may actually be important to the performance of certain
applications. It is worth considering whether other optimizations could be used
for determining equality that improve performance in these cases. This is also
left for future work.

\section{Putting the pieces together}
\label{sec:numberType}

With the definitions and properties of \code{AdjoinSqrt} and \code{Filtered}
above, we can put all of the pieces together to form the example number type
described at the beginning of the chapter and conclude that it is an ordered
field.

\begin{definition}[{Approximation for $A[\sqrt{n}]$}]
  \label{def:adjoinApprox}
  \uses{def:adjoin, def:approxStrategy}
  \lean{AdjoinSqrt.toReal, CGLean.ApproxSqrt, AdjoinSqrt.approxStrategy}
  \leanok
  Given an approximation strategy $f : A → α$ with embedding $i : A → ℝ$, an
  element $n ∈ A$ with $i(n) \geq 0$, and an element $\hat{S} ∈ α$ with
  $\sqrt{i(n)} ∈ \hat{S}$, we define an approximation strategy for
  $A[\sqrt{n}]$ with embedding
  $$i_{\sqrt{n}}(a_1 + a_n\sqrt{n}) := i(a_1) + i(a_n)\sqrt{i(n)}$$ and
  approximation $$f_{\sqrt{n}}(a_1 + a_n\sqrt{n}) := f(a_1) + f(a_n)\hat{S}$$
\end{definition}

\begin{lemma}[$f_{\sqrt{n}}$ is a valid approximation strategy]
  \label{thm:adjoinApproxValid}
  \uses{def:adjoinApprox, thm:adjoinOrderedRing}
  \lean{AdjoinSqrt.toRealOrderHom, AdjoinSqrt.toReal_nonneg, AdjoinSqrt.toReal_monotone, AdjoinSqrt.approxStrategy}
  \leanok
  If $A$ is a linearly ordered field with a positive non-square $n ∈ A$, then
  $i_{\sqrt{n}}$ is an embedding and $f_{\sqrt{n}}$ is a valid
  approximation strategy.
\end{lemma}

\begin{proof}
  \leanok
  We need to check that $i_{\sqrt{n}} : A[\sqrt{n}] → ℝ$ is an order embedding
  and that $f_{\sqrt{n}} : A[\sqrt{n}] → α$ is a valid approximation strategy.

  Most of the requirements for $i_{\sqrt{n}}$ are immediate from the definitions. To check
  the ordering property of $i_{\sqrt{n}}$, we first prove that if $0 \leq a$ then $0 \leq
  i_{\sqrt{n}}(a)$, from which the main result follows.

  To check that, we consider the cases where $\norm{a} \geq 0$ and $\norm{a}
  \leq 0$ separately. In each case, one of the two terms in $i_{\sqrt{n}}(a)$
  is nonnegative by assumption, and by squaring you can check that the nonnegative
  one dominates.

  Now we need to show that $f_{\sqrt{n}} : A[\sqrt{n}] → α$ is a valid strategy, i.e.
  that $i_{\sqrt{n}}(a) ∈ f_{\sqrt{n}}(a)$ for all $a ∈ A[\sqrt{n}]$. We have
  $i_{\sqrt{n}}(a_1 + a_n\sqrt{n}) = i(a_1) + i(a_n)\sqrt{i(n)}$. Since $f$ is a valid
  strategy we have $i(a_1) ∈ f(a_1)$ and $i(a_n) ∈ f(a_n)$, and by assumption
  $\sqrt{i(n)} ∈ \hat{S}$; since $α$ is an approximation type, we can put these
  together to conclude $i(a_1) + i(a_n)\sqrt{i(n)} ∈ f(a_1) + f(a_n)\hat{S} =
  f_{\sqrt{n}}(a_1 + a_n\sqrt{n})$ as required.
\end{proof}

\begin{corollary}[Example number type]
  \label{thm:numberType}
  \uses{def:adjoin, def:filtered, def:dyadicInterval}
  \lean{CGLean.NumberType.isStrictOrderedRing, CGLean.NumberType.instFieldNum, CGLean.NumberType.instLinearOrderNum}
  \leanok
  \code{Filtered $ℚ[\sqrt{2}][\sqrt{3}]$} is a linearly ordered field.
\end{corollary}
\begin{proof}
  \uses{thm:adjoinOrderedField, thm:adjoinApproxValid, thm:filteredField, thm:filteredOrdered, thm:dyadicApprox}
  \leanok
  $ℚ$ is clearly a linearly ordered field. The fact that $2$ is a positive
  non-square in $ℚ$ is well-known, and it can be used to prove that $3$ is also
  a positive non-square in $ℚ[\sqrt{2}]$ (the proof by contradiction is easy).
  Therefore $ℚ[\sqrt{2}][\sqrt{3}]$ is a linearly ordered field.

  Applying Lemma~\ref{thm:adjoinApproxValid} twice, starting from the
  rationals approximated by dyadic intervals, gives an approximation strategy
  for $ℚ[\sqrt{2}][\sqrt{3}]$. Plugging this into the construction of
  \code{Filtered} with our dyadic interval approximation gives a linearly
  ordered field structure on \code{Filtered $ℚ[\sqrt{2}][\sqrt{3}]$} as
  required.
\end{proof}

% vim: ts=2 sw=2 et ai
